% !TeX root = ../../Book3.tex
% 纲要标题：### 16.2 Ext 刻画
% 纲要位置：计划纲要.md，第 2466 行。

\section{Ext 刻画}
\label{b3:sec:1602}

一个正则元素在模上作用为单射，却把剩余域零化。把这两种作用放进 Hom 的消解计算，短正合列便使首次非零的上同调次数改变一。正则序列的长度因此有另一种表达：从剩余域到模的 Ext 首次非零次数。局部化时，剩余域随素理想改变，还需要把这种计算与支集维数联系起来。


% 纲要内容（仅作写作计划，尚非教材正文）：
%
% * Ext 的两变量长正合列、零化作用、有限性与局部化
% * 以剩余域的 Ext 检测深度
% * 与正则序列定义的等价
% * Ischebeck 消失、关联素理想维数界与深度的局部化估计
% * 正则序列、Koszul 同调与 Ext 的深度刻画对照
%

\subsection{剩余域的 Ext 与深度检测}
\label{b3:subsec:1602:01}

\begin{definition}\label{b3:def:ext-minimum-background}
设 \(R\) 为交换环，\(N,M\) 为 \(R\) 模，选自由消解 \(P_\bullet\rightarrow N\rightarrow0\)，其微分记为 \(\mathrm{d}_i:P_i\rightarrow P_{i-1}\)（\(i\ge1\)）。取上链复形
\[
C^i=\operatorname{Hom}_R(P_i,M)\quad(i\ge0),\qquad
\delta^i(f)=f\circ \mathrm{d}_{i+1}.
\]
令 \(C^i=0\) 对 \(i<0\) 成立，\(\delta^{-1}=0\)。因 \(\mathrm{d}_i\mathrm{d}_{i+1}=0\)，有 \(\delta^{i+1}\delta^i=0\)。定义
\[
\operatorname{Ext}_R^i(N,M)=H^i(C^\bullet)
=\ker\delta^i/\operatorname{im}\delta^{i-1}\quad(i\ge0),
\]
并令负次数的 Ext 为零。
\end{definition}
\symidx{\(\operatorname{Ext}_R^i(N,M)\)：Ext 模}

由\cref{b3:lem:free-resolution-existence}，每个模都有自由消解。\cref{b3:lem:free-resolution-comparison}给出不同消解之间的比较映射及其链同伦唯一性；施加 \(\operatorname{Hom}_R(-,M)\) 后，同伦仍给出相同的上同调映射。因此 Ext 不依赖选择，且次数零自然同构于 \(\operatorname{Hom}_R(N,M)\)。同态 \(N'\rightarrow N\) 提升为消解的链映射后，经 Hom 反向诱导 Ext 映射；同态 \(M\rightarrow M'\) 则经后复合正向诱导映射。比较引理使这些映射与选择无关，并保持恒等映射及复合，所以 Ext 在第一变量反变、第二变量协变。

\begin{lemma}[Ext 的长正合列与零化作用]\label{b3:lem:ext-basic-properties}
设 \(R\) 为交换环，\(N,M\) 为 \(R\) 模。
\begin{enumerate}
\item 短正合列在 Ext 的第二变量给出协变长正合列，在第一变量给出反变长正合列。
\item 若 \(a\in R\) 满足 \(aN=0\) 或 \(aM=0\)，则 \(a\operatorname{Ext}_R^i(N,M)=0\) 对所有 \(i\ge0\) 成立。
\end{enumerate}
\end{lemma}
\begin{proof}
第二变量的短正合列经每个 \(\operatorname{Hom}_R(P_i,-)\) 仍正合，因为自由模的映射可对基逐个提升。于是得到短正合的上链复形。其连接映射这样构造：提升商复形中的一个圈，取其微分，结果落在子复形中且为圈；改变提升或原圈代表只改变一个边界。这也给出正合性：连接类为零恰好表示可以修正提升使其成为圈，其他两处分别对应圈来自子复形和边界能够提升。

第一变量的短正合列记为 \(0\rightarrow N'\rightarrow N\rightarrow N''\rightarrow0\)，给两端选自由消解 \(P',P''\)。令 \(Z'_{-1}=N'\)、\(Z_{-1}=N\)、\(Z''_{-1}=N''\)。递归构造 \(P\) 时，假设已有
\[
0\longrightarrow Z'_{i-1}\xrightarrow{j}Z_{i-1}
\xrightarrow{q}Z''_{i-1}\longrightarrow0.
\]
两端消解给出满射 \(\rho'_i:P'_i\rightarrow Z'_{i-1}\)、\(\rho''_i:P''_i\rightarrow Z''_{i-1}\)，记其核为 \(Z'_i,Z''_i\)。因 \(P''_i\) 自由，可取提升 \(\ell_i:P''_i\rightarrow Z_{i-1}\)，使 \(q\ell_i=\rho''_i\)。定义
\[
P_i=P'_i\oplus P''_i,\qquad
\rho_i(u,v)=j\rho'_i(u)+\ell_i(v).
\]
它为满射：先从 \(P''_i\) 提升 \(qz\)，余下的差属于 \(jZ'_{i-1}\)，再从 \(P'_i\) 提升该差。令 \(Z_i=\ker\rho_i\)。包含第一分量与投影第二分量使
\[
0\longrightarrow Z'_i\longrightarrow Z_i\longrightarrow Z''_i\longrightarrow0
\]
正合。当前这一层的构造放在三列图中如下：中排是刚选出的自由模，向下满射到上一层的三个核，向上则是本层新得到的三个核。各行均短正合，中排还分裂。
\[
\begin{tikzcd}[column sep=2em,row sep=2.5em]
0\arrow[r]&Z'_i\arrow[r,hook]\arrow[d,hook]
 &Z_i\arrow[r,two heads]\arrow[d,hook]
 &Z''_i\arrow[r]\arrow[d,hook]&0\\
0\arrow[r]&P'_i\arrow[r,hook]\arrow[d,two heads,"\rho'_i"']
 &P'_i\oplus P''_i\arrow[r,two heads]\arrow[d,two heads,"\rho_i"]
 &P''_i\arrow[r]\arrow[d,two heads,"\rho''_i"]&0\\
0\arrow[r]&Z'_{i-1}\arrow[r,hook,"j"']
 &Z_{i-1}\arrow[r,two heads,"q"']
 &Z''_{i-1}\arrow[r]&0
\end{tikzcd}
\]
构造 \(\rho_i\) 时，先用 \(\ell_i\) 提升右端，左端再修正误差。尤其，对 \(v\in Z''_i\)，\(\ell_i(v)\) 落在 \(jZ'_{i-1}\) 中，可取 \(u\) 使 \(j\rho'_i(u)=-\ell_i(v)\)，故 \((u,v)\in Z_i\) 提升 \(v\)；投影的核恰是 \(Z'_i\)。上排因此具有与下排相同的形式，可以作为下一层的输入。继续递归，取 \(\rho_0\) 为增广映射，并在 \(i\ge1\) 时把 \(\rho_i\) 与 \(Z_{i-1}\hookrightarrow P_{i-1}\) 复合，得到中间的自由消解 \(P\)。其包含与投影保持微分，组成逐项分裂的短正合消解
\[
0\longrightarrow P'\longrightarrow P\longrightarrow P''\longrightarrow0,
\qquad P_i=P'_i\oplus P''_i.
\]
于是每个次数有分裂短正合列
\[
0\longrightarrow\operatorname{Hom}_R(P''_i,M)
\longrightarrow\operatorname{Hom}_R(P_i,M)
\longrightarrow\operatorname{Hom}_R(P'_i,M)\longrightarrow0.
\]
前段圈提升的构造就给出第一变量的反变长正合列。这一逐项分裂的消解构造称为马蹄构造；连接映射与比较映射相容，因而所得长正合列自然。

若 \(aN=0\)，则 \(a\operatorname{id}_{P_\bullet}\) 与零映射都提升 \(N\) 上的零同态。\cref{b3:lem:free-resolution-comparison}给出链同伦 \(a\operatorname{id}=\mathrm{d}h+h\mathrm{d}\)。施加 Hom 后，乘法 \(a\) 在上同调上为零。若 \(aM=0\)，Hom 复形各项上的乘法 \(a\) 已经为零。
\end{proof}

\begin{proposition}[有限生成的 Ext 与局部化]\label{b3:prop:finite-ext-localization}
设 \(R\) 为交换 Noether 环，\(N,M\) 为有限生成的 \(R\) 模，\(i\ge0\) 为整数。则 \(\operatorname{Ext}_R^i(N,M)\) 是有限生成的 \(R\) 模，并且任意乘法集 \(S\subseteq R\) 都给出自然同构
\[
S^{-1}\operatorname{Ext}_R^i(N,M)
\cong\operatorname{Ext}_{S^{-1}R}^i(S^{-1}N,S^{-1}M).
\]
\end{proposition}
\begin{proof}
因 \(N\) 有限生成，可先选有限秩自由模到 \(N\) 的满射。每个核都由 Noether 性有限生成，所以可以逐次选有限秩自由模满射，得到各项均为有限秩的自由消解 \(P\)。每个 \(\operatorname{Hom}_R(P_i,M)\) 是有限生成模，其核和像也有限生成，故上同调有限生成。对有限秩自由模有
\[
S^{-1}\operatorname{Hom}_R(P_i,M)
\cong\operatorname{Hom}_{S^{-1}R}(S^{-1}P_i,S^{-1}M),
\]
且这些同构保持微分。局部化又保持消解的正合性、核及像，因此取上同调得到所述自然同构。
\end{proof}

\subsection{Ext 刻画与正则序列定义的等价性}
\label{b3:subsec:1602:02}

\begin{theorem}\label{b3:thm:depth-ext}
设 \((R,\mathfrak m)\) 为交换 Noether 局部环，\(k=R/\mathfrak m\) 为其剩余域，\(M\ne0\) 为有限生成的 \(R\) 模。则
\[
\operatorname{depth}_RM
=\inf\{\,i\ge0\mid\operatorname{Ext}_R^i(k,M)\ne0\,\}.
\]
特别地，右侧集合非空，其最小值等于正则序列定义的深度。
\end{theorem}
\begin{proof}
对 \(t=\operatorname{depth}_RM\) 归纳。\(t=0\) 时，由\cref{b3:cor:depth-zero-socle}，\(\operatorname{Ext}_R^0(k,M)=\operatorname{Hom}_R(k,M)\ne0\)。设 \(t>0\)，选 \(x\in\mathfrak m\) 为 \(M\) 非零因子，令 \(\bar M=M/xM\)。由\cref{b3:cor:regular-quotient-depth}，\(\operatorname{depth}_R\bar M=t-1\)。

短正合列 \(0\rightarrow M\xrightarrow{x}M\rightarrow\bar M\rightarrow0\) 给出 Ext 长正合列。因为 \(xk=0\)，\cref{b3:lem:ext-basic-properties}说明所有 Ext 上的乘法 \(x\) 为零，故每个 \(i\ge0\) 有
\[
0\longrightarrow\operatorname{Ext}_R^i(k,M)
\longrightarrow\operatorname{Ext}_R^i(k,\bar M)
\longrightarrow\operatorname{Ext}_R^{i+1}(k,M)\longrightarrow0.
\]
又因 \(x\) 在 \(M\) 上单射，\(\operatorname{Hom}_R(k,M)=0\)。归纳假设说中间项在 \(i<t-1\) 为零、在 \(i=t-1\) 非零。上述短正合列遂给出 \(\operatorname{Ext}_R^i(k,M)=0\) 对 \(i<t\) 成立，并给出
\[
\operatorname{Ext}_R^t(k,M)\cong\operatorname{Ext}_R^{t-1}(k,\bar M)\ne0.
\]
这完成归纳。
\end{proof}

例如在 \(R=k[X,Y]_{(X,Y)}\) 上，用 \(X,Y\) 的 Koszul 消解计算 \(\operatorname{Ext}_R^i(k,R)\)。对偶复形的两个矩阵是 \(\binom XY\) 与 \((-Y\ \ X)\)，其前两处上同调为零，最后的余核为 \(k\)。所以深度为二。在 \(R=k[\varepsilon]/(\varepsilon^2)\) 上，映射 \(k\rightarrow R\)、\(1\mapsto\varepsilon\) 已在次数零非零，故深度为零，不必先构造整个无限消解。

\ProseParagraphBreak
在第一个例子中，若交换两个变量，\(R\) 自身已是自由模，可以用只含零次项的消解计算，故 \(\operatorname{Ext}_R^{i>0}(R,k)=0\)。这与 \(\operatorname{Ext}_R^2(k,R)\cong k\) 同时成立：深度判据把 \(k\) 放在第一变量、被检测的模放在第二变量。后面的投射维数判据检测的是被消解的第一变量。Tor 的变量对称性不能照搬给 Ext。

\subsection{Ischebeck 消失与深度的局部化估计}
\label{b3:subsec:1602:03}

设 \((R,\mathfrak m)\) 为交换 Noether 局部环，\(k=R/\mathfrak m\)。局部化 Ext 的公式不能把闭点的剩余域直接变成其他素理想处的剩余域。若 \(\mathfrak p\subsetneq\mathfrak m\)，则 \(k_{\mathfrak p}=0\)，而 \(\kappa(\mathfrak p)\) 非零。比较两个地方的深度，需要额外的维数估计。Ischebeck 的消失定理把第一变量的支集维数与第二变量的深度联系起来。\footnote{弗里德里希·格哈特·伊舍贝克（Friedrich Gerhart Ischebeck，1940--2022），德国数学家，主要研究交换代数和同调代数，著有《Einladung zur Zahlentheorie》。}

\begin{lemma}\label{b3:lem:ischebeck-vanishing}
设 \((R,\mathfrak m)\) 为交换 Noether 局部环，\(M,N\ne0\) 为有限生成的 \(R\) 模，\(t=\operatorname{depth}_RM\)、\(s=\dim_RN\)。则
\[
\operatorname{Ext}_R^i(N,M)=0\qquad(0\le i<t-s).
\]
\end{lemma}
\begin{proof}
对 \(s\) 归纳，先说明如何归约到素循环模。任一非零有限生成模都有有限素因子滤过：在非零商中选一个关联素理想，便得到同构于 \(R/\mathfrak q\) 的子模，再在余下的商中重复；相应的子模严格升链由 Noether 性终止。将所得滤过记为
\[
0=N_\ell\subset N_{\ell-1}\subset\cdots\subset N_0=N,
\qquad N_j/N_{j+1}\cong R/\mathfrak q_j.
\]
各因子的支集包含于 \(\operatorname{Supp}N\)，所以 \(\dim(R/\mathfrak q_j)\le s\)。若对某个次数 \(i\) 所有因子的 Ext 都为零，从 \(N_\ell=0\) 开始对滤过长度归纳。短正合列
\(0\rightarrow N_{j+1}\rightarrow N_j\rightarrow R/\mathfrak q_j\rightarrow0\)
的第一变量长正合列给出
\[
\operatorname{Ext}_R^i(R/\mathfrak q_j,M)
\longrightarrow\operatorname{Ext}_R^i(N_j,M)
\longrightarrow\operatorname{Ext}_R^i(N_{j+1},M).
\]
外两项为零便使中项为零，逐层推出 \(\operatorname{Ext}_R^i(N,M)=0\)。因此，只需对 \(N=R/\mathfrak q\) 证明指定范围内的消失。

当 \(s=0\) 时，各因子的素理想只能为 \(\mathfrak m\)，因子就是 \(k\)。\cref{b3:thm:depth-ext}给出各因子的 Ext 在 \(i<t\) 时为零，前段长度归纳完成基例。设 \(s>0\)。若 \(\dim(R/\mathfrak q)<s\)，由维数归纳假设得到消失；其消失范围至少包含 \(i<t-s\)。剩下 \(\dim(R/\mathfrak q)=s\) 的情形。选 \(x\in\mathfrak m\setminus\mathfrak q\)，则乘法 \(x\) 在 \(R/\mathfrak q\) 上单射，有
\[
0\longrightarrow R/\mathfrak q\xrightarrow{x}R/\mathfrak q
\longrightarrow R/(\mathfrak q,x)\longrightarrow0.
\]
由\cref{b3:prop:regular-quotient-dimension}，末项维数为 \(s-1\)。归纳假设使其 Ext 在次数小于 \(t-s+1\) 时为零。对 \(0\le i<t-s\)，第一变量长正合列中有
\[
\operatorname{Ext}_R^i(R/\mathfrak q,M)
\xrightarrow{x}\operatorname{Ext}_R^i(R/\mathfrak q,M)
\longrightarrow\operatorname{Ext}_R^{i+1}(R/(\mathfrak q,x),M)=0.
\]
因此乘法 \(x\) 为满射。\cref{b3:prop:finite-ext-localization}说明该 Ext 模有限生成；因 \(x\in\mathfrak m\)，Nakayama 引理使其为零。再按素因子滤过的长度归纳，便得到原模 \(N\) 的结论。
\end{proof}

\begin{corollary}\label{b3:cor:associated-depth-bound}
设 \((R,\mathfrak m)\) 为交换 Noether 局部环，\(M\ne0\) 为有限生成的 \(R\) 模。若 \(\mathfrak q\in\operatorname{Ass}_RM\)，则
\[
\operatorname{depth}_RM\le\dim(R/\mathfrak q).
\]
\end{corollary}
\begin{proof}
关联素理想给出非零单射 \(R/\mathfrak q\rightarrow M\)。若维数小于深度，上述引理将使 \(\operatorname{Hom}_R(R/\mathfrak q,M)=0\)，与该单射矛盾。
\end{proof}

\begin{theorem}\label{b3:thm:depth-localization}
设 \((R,\mathfrak m)\) 为交换 Noether 局部环，\(M\ne0\) 为有限生成的 \(R\) 模，\(\mathfrak p\in\operatorname{Supp}_RM\)。则
\[
\operatorname{depth}_RM
\le\operatorname{depth}_{R_{\mathfrak p}}M_{\mathfrak p}
+\dim(R/\mathfrak p).
\]
\end{theorem}
\begin{proof}
在 \(\mathfrak p\) 中连续选取 \(M\) 正则元素，直至不能延长，记所得序列长度为 \(r\)，末端商为 \(L\)。每加进一个正则元素，子模 \((x_1,\ldots,x_j)M\) 严格增大，所以 \(M\) 的 Noether 性保证过程终止。\(\mathfrak p\) 中所有元素都是 \(L\) 的零因子；而 \(\operatorname{Ass}_RL\) 有限，故由有限素理想回避，存在 \(\mathfrak q\in\operatorname{Ass}_RL\)，使 \(\mathfrak p\subseteq\mathfrak q\)。\cref{b3:cor:regular-quotient-depth}和\cref{b3:cor:associated-depth-bound}给出
\[
\operatorname{depth}_RM=r+\operatorname{depth}_RL
\le r+\dim(R/\mathfrak q)\le r+\dim(R/\mathfrak p).
\]
原序列局部化后仍逐次单射，而且所有元素都在 \(\mathfrak pR_{\mathfrak p}\) 中。\(M_{\mathfrak p}\ne0\) 且有限生成，Nakayama 引理保证每个商非零，所以这是一个长度 \(r\) 的 \(M_{\mathfrak p}\) 正则序列。于是 \(r\le\operatorname{depth}_{R_{\mathfrak p}}M_{\mathfrak p}\)，代入便得结论。
\end{proof}

有限生成条件使关联素理想集及 Ext 模受 Noether 性控制，也使 Nakayama 引理适用。例如，设 \(R\) 为离散赋值环，\(\pi\) 为其统一化元，\(K\) 为分式域。乘法 \(\pi\) 在 \(K\) 上既单射又满射，因而 \(K/\pi K=0\)；它不满足正则序列的末端非零条件。

\ProseParagraphBreak
\cref{b3:thm:maximal-regular-length,b3:thm:depth-ext}将正则序列的长度与两种同调计算联系起来，\cref{b3:cor:depth-zero-socle}给出深度零的等价条件。这些刻画集中列于\cref{b3:tab:depth-characterizations}。

\begin{center}
\begin{minipage}{\textwidth}
\makeatletter
\def\@captype{table}
\makeatother
\centering
\caption[深度的三种刻画]{深度的三种刻画。设 \((R,\mathfrak m)\) 为交换 Noether 局部环，\(M\ne0\) 为有限生成的 \(R\) 模，\(k=R/\mathfrak m\)，\(y=(y_1,\ldots,y_e)\) 为 \(\mathfrak m\) 的任意有限生成组。前三行给出同一个数值，末行列出该数值为零时的等价条件。}
\label{b3:tab:depth-characterizations}
\begin{tabularx}{\textwidth}{@{}>{\raggedright\arraybackslash}p{0.25\textwidth}>{\raggedright\arraybackslash}X@{}}
\toprule
刻画 & 数值或条件\\
\midrule
正则序列 & \(\mathfrak m\) 中任一极大 \(M\) 正则序列的长度\\
Koszul 同调 & \(e-\max\{\,i\mid H_i(y;M)\ne0\,\}\)\\
Ext & \(\min\{\,i\ge0\mid\operatorname{Ext}_R^i(k,M)\ne0\,\}\)\\
深度零 & \(\mathfrak m\in\operatorname{Ass}_RM\ \Longleftrightarrow\ \operatorname{Soc}_RM\ne0\)\\
\bottomrule
\end{tabularx}
\end{minipage}
\end{center}
